\section{결론}
감사된 시뮬레이션 결과, 제어기 순위는 목적함수, 모델 불일치, 누락 동특성에 따라 달라졌다. 동일 예산 튜닝은 공칭 순위를 LQI--PID--LQR에서 LQR--LQI--PID로 바꾸지만 작은 센서 지연은 튜닝 LQR의 이점을 제거한다. 적분 동작의 가장 명확한 이점은 일정 부하 제거이다. 고정 시간 RMSE에는 같은 백분율의 열전달 불확실성보다 열용량 불확실성이 더 큰 영향을 미쳤다. 포화 처리 수정으로 기존 부정적 결과가 무효화되었고 잡음 차이는 통계적으로 명확하지만 실용적으로 작다. 구간 분석은 선언한 연속 매개변수 상자에서 기존 LQR과 LQI의 비포화 강인 안정성을 증명하지만, PID와 모든 포화 폐루프의 근거는 명시적으로 제한된다. 저장소는 새로운 제어 알고리즘이나 실제 공정 검증이 아니라 재현 가능한 연구 결과를 제공한다. 향후 연구는 물리 실험을 요구하지 않고 구간 분석, 지연 강인 튜닝, 유색 잡음, 고차 열 네트워크로 확장할 수 있다.

\section*{재현성 선언}
실행 가능한 모든 연구 코드는 Rust로 작성되었다. 저장소는 Rust 1.100.0과 의존성 버전을 고정하고, 설정 및 시드를 기록하며, 계산 결과에서 표와 벡터 그림을 생성하고, 제공된 템플릿 절차로 논문을 다시 빌드한다. 전체 연구는 \texttt{cargo run --release -- run-all} 실행 후 \texttt{paper/}에서 \texttt{make}로 재현한다.
